% !TEX root = ../00_Main/Main.tex
\documentclass[../00_Main/Main.tex]{subfiles}
\begin{document}

\label{sec:discussion}

Los resultados respaldan la hip\'otesis H1 y no respaldan la H2: s\'olo
una de las seis reglas, el ensamble ponderado, supera al Transformer, con
una diferencia de tama\~no despreciable (Secci\'on~\ref{sec:ens-best}).

\subsection{Interpretaci\'on de la ventaja del Transformer}\label{sec:disc-individual}

La ventaja del Transformer no puede atribuirse a la ventana de estados: el
DQN recurrente recibe la misma ventana de seis estados y en
generalizaci\'on queda un $4\%$ por debajo del Transformer e incluso un
$1\%$ por debajo del DQN, que s\'olo observa el estado actual. Tampoco se
debe al tama\~no, porque el Transformer tiene un $21\%$ menos de
par\'ametros que el DQN recurrente. Descartados esos dos
factores, la mejora se atribuye a la atenci\'on causal, que pondera cada
estado de la ventana seg\'un su relaci\'on con el estado actual en lugar
de resumirla en un \'unico estado oculto, como hace la LSTM. La ventaja
se conserva en generalizaci\'on: el Transformer obtiene la mayor media en
el $73\%$ de los escenarios ($16$ de $22$;
Figura~\ref{fig:heatmap-global}) y su media all\'i supera en un $0{,}3\%$
a la de la referencia. Su mayor p\'erdida ocurre con secuencias m\'as largas
que las de entrenamiento, donde cae un $7\%$ y queda un $3\%$ por debajo
del DQN recurrente. Como no se evalu\'o una variante sin atenci\'on, su
efecto no puede separarse del de la normalizaci\'on previa a cada subcapa~\citep{Parisotto2020} ni del de
los hiperpar\'ametros.

\subsection{Efecto de la regla de combinaci\'on}\label{sec:ens-best}

Los seis ensambles combinan los mismos modelos y ponderan a cada miembro
por su confianza; s\'olo cambia la regla (Tabla~\ref{tab:ens-freq}):

\begin{itemize}
    \item En el \textbf{ensamble ponderado} la decisi\'on no la determina
    el Transformer, pese a que tiene el $82\%$ del peso base: $\tau = 15$ aplana las distribuciones
    de los miembros y el \emph{prior} cambia la acci\'on votada en el
    $41\%$ de las decisiones de la referencia y el $33\%$ de las de
    generalizaci\'on. Su ventaja se asocia a la combinaci\'on de una
    preferencia d\'ebil de los modelos con una regla fija que asigna
    aproximadamente seg\'un la severidad.
    \item La \textbf{compuerta del Transformer} lo sigue en el $97\%$ y el
    $95\%$ de las decisiones, y por eso alcanza un desempe\~no
    equivalente.
    \item En el \textbf{voto suave} y el \textbf{voto por acci\'on} los
    pesos base favorecen al DQN recurrente sobre el Transformer; en la
    extrapolaci\'on de severidad, la ponderaci\'on por confianza no
    produjo que predominara el miembro con mayor desempe\~no individual.
    \item El \textbf{consenso con \emph{prior}} mejora la referencia y el
    peor escenario del consenso, aunque su \emph{prior} s\'olo cambia el
    $2\%$ y el $4\%$ de las decisiones; la mejora no puede atribuirse al
    \emph{prior}.
\end{itemize}

Los datos son compatibles con que la ventaja del ensamble ponderado
provenga del \emph{prior} y de la temperatura alta m\'as que de la
ponderaci\'on por confianza; con un \'unico ensamble que supera al Transformer no puede
separarse el aporte de cada componente. Los datos s\'i muestran que esa
ponderaci\'on, por s\'i sola, no basta para superar al mejor miembro: los
dos ensambles que s\'olo promedian o votan con ella quedan por debajo del
Transformer individual.

\biblio
\end{document}
